\chapter{Best practices}
\label{ch:bestpractices}

Concentrated lessons from real-world deployments. Every section is
optional — pick what fits your scale.

\section{Storage hierarchy}

\begin{figure}[ht]
  \centering
  \includegraphics[width=0.95\textwidth]{locations.png}
  \caption{The Locations page on a live install: a household
    hierarchy with a study room, an in-room bookcase, ten shelves,
    and two flat-level rooms — each node carrying its L-code and a
    live volume count.}
  \label{fig:locations}
\end{figure}

\begin{itemize}
  \item \textbf{Three levels are usually enough.} Room → shelf →
        row covers most households. A fourth level (e.g.\ ``box
        inside row'') becomes paperwork you have to maintain.
  \item \textbf{Name locations the way you talk about them.}
        ``Living-room shelf, top'' beats ``LR-S01-R01''. The names
        are visible on barcode labels; treat them as user-facing.
  \item \textbf{Print L-codes once.} Re-printing because you
        rebalanced shelves is wasted effort — L-codes are stable
        identifiers; the human-readable name on the label can drift
        and that is fine.
\end{itemize}

\section{V-code rolls}

\begin{itemize}
  \item Print V-codes in rolls of 50 to 100. Pre-printed stickers in
        a holder beside the scanner station make cataloging
        \emph{fast}.
  \item \textbf{Stick the V-code before scanning.} It removes a
        future failure mode where you mislabel a book and have to
        scan-find it again.
  \item \textbf{Place the sticker consistently.} Inside back cover,
        upper-right corner, parallel to the spine — pick a rule
        and stick to it. Future-you doing an audit will thank
        present-you.
\end{itemize}

\section{Genre versus Dewey}

mybibli supports both broad genres and Dewey codes\index{Dewey}.
Pick one as your primary navigation aid:

\begin{itemize}
  \item \textbf{Use genres} when the audience browses by mood
        (``I want a thriller'') and the collection is small enough
        that ``Thriller'' is granular enough.
  \item \textbf{Use Dewey} when the audience browses by topic
        (``books on 18th-century French history'') or when the
        collection has reached the size where ``Non-fiction'' is no
        longer a useful filter.
  \item Using both is fine; mybibli does not force a choice.
\end{itemize}

\section{Audit cadence}

\begin{itemize}
  \item \textbf{Quarterly} for collections under a few hundred
        volumes.
  \item \textbf{Monthly} for collections above that, especially if
        more than one person catalogs.
  \item Always run an audit \emph{after} a major shelf rebalance —
        moving books physically without telling mybibli is the
        single biggest source of ``where is this volume?'' bugs.
\end{itemize}

\section{User accounts}

\begin{itemize}
  \item \textbf{One human, one account.} Sharing a single
        \texttt{admin} account between two people loses the audit
        trail and makes the last-active-admin guard less useful.
  \item \textbf{Most editors should be Librarians, not Admins.} The
        admin panel handles installation-time concerns and rare
        cleanups; you do not need admin rights to catalog or loan.
  \item \textbf{Rotate the seeded admin password on first login.}
        On 1.1.0+ the setup wizard forces this; on legacy 1.0.x
        installs do it manually before exposing the instance to a
        network.
\end{itemize}

\section{Loans}

\begin{itemize}
  \item \textbf{Keep borrowers identifiable.} ``Marie L.'' beats
        ``Marie'' the moment the second Marie shows up. Email or
        phone in the notes field helps for chase-ups.
  \item \textbf{Run the overdue list weekly.} The dashboard
        indicator nags you for a reason — overdue loans
        compounded over a year become books you no longer
        own.
  \item \textbf{Mark damage at return time.} Switching a returned
        volume to \emph{Worn} or \emph{Damaged} captures the
        information while the borrower is in front of you; chasing
        the cost later from memory rarely works.
\end{itemize}

\section{Metadata providers}

\begin{itemize}
  \item \textbf{Keep keys minimal.} You do not need all three
        keyed providers from day one — BnF + Open Library cover
        books for free. Add OMDb + TMDb the day you catalog
        your first DVD.
  \item \textbf{Treat keys as secrets.} \texttt{.env} is in
        \texttt{.gitignore}; keep backups of \texttt{.env}
        encrypted (chapter~\ref{ch:backup}).
\end{itemize}

\section{Upgrades}

\begin{itemize}
  \item \textbf{Read release notes before pulling \texttt{:latest}.}
        See chapter~\ref{ch:upgrade}.
  \item \textbf{Back up immediately before any upgrade.} Even
        ``additive'' migrations have non-zero failure surface;
        a same-minute backup turns a worry into a chore.
  \item \textbf{Pin a version tag for production.} Auto-updates
        for a household instance are fine; for a community library,
        an unattended major bump can interrupt service.
\end{itemize}
